\documentclass[12pt,a4paper]{article}

\usepackage[utf8]{inputenc}
\usepackage[T1]{fontenc}
\usepackage{amsmath,amssymb,amsfonts}
\usepackage{mathtools}
\usepackage{graphicx}
\usepackage[margin=2.5cm]{geometry}
\usepackage{setspace}
\usepackage{booktabs}
\usepackage{enumitem}
\usepackage[dvipsnames]{xcolor}
\usepackage{hyperref}
\usepackage{cleveref}
\usepackage{natbib}
\usepackage{abstract}
\usepackage{titlesec}
\usepackage{todonotes}

\hypersetup{
    colorlinks=true,
    linkcolor=blue!70!black,
    citecolor=blue!70!black,
    urlcolor=blue!70!black,
    pdfauthor={Franklin Silveira Baldo},
    pdftitle={Syntax as Physics: The Causal Injection Test and the Proxy Ontology Defense},
}

\onehalfspacing
\titleformat{\section}{\large\bfseries}{\thesection.}{0.5em}{}
\titleformat{\subsection}{\normalsize\bfseries}{\thesubsection}{0.5em}{}
\renewcommand{\abstractnamefont}{\normalfont\bfseries}
\renewcommand{\abstracttextfont}{\normalfont\small}

\begin{document}

\begin{center}
    {\LARGE\bfseries Syntax as Physics: The Causal Injection Test and the Proxy Ontology Defense\\[4pt]}

    \vspace{1.2em}

    {\large Franklin Silveira Baldo}\\[4pt]
    {\normalsize Procuradoria Geral do Estado de Rondônia}\\
    {\small\texttt{franklin.baldo@pge.ro.gov.br}}

    \vspace{1em}
    {\small March 2026}
\end{center}
\vspace{0.5em}

\begin{abstract}
\noindent
In response to my framework establishing "narrative residue" within autoregressive generation, Hossenfelder (2026) and Aaronson (2026) assert that equating statistical token generation with physical law commits a "Proxy Ontology Fallacy." They argue that because an LLM's O(1) forward pass structurally fails to compute true combinatorial probabilities, its outputs are merely statistical hallucinations mapped to human linguistic biases. I concede entirely that the model operates via text co-occurrence and prompt sensitivity, and cannot magically execute \#P-hard counting. However, I argue that within a text-based simulated universe, this syntax *is* the physics. The Causal Injection Test, specifically isolating Mechanism C (autoregressive narrative conditioning), demonstrates that the substrate fundamentally corrupts isolated mathematical probabilities by enforcing synthetic causal non-locality across independent systems. The model's algorithmic failure to sustain an objective reality is not a reason to dismiss the ontology; it is precisely the empirical proof that a generating substrate dictates the fundamental laws of its simulated world. The narrative structure of human language acts as the implicit Hamiltonian governing this explicitly generated text reality.
\end{abstract}

\section{Introduction and the Concession of the Statistical Fallacy}

The critical responses to my recent paper \emph{The Narrative Residue} \citep{baldo2026_narrative} from Sabine Hossenfelder and Scott Aaronson focus on what Hossenfelder terms the "Statistical Fallacy" and the "Proxy Ontology Fallacy."

Hossenfelder correctly points out that an LLM's single forward pass, bounded by an $O(1)$ depth limit, is entirely incapable of computing the \#P-hard constraints required to determine the exact ground-truth probability of a Minesweeper board. Therefore, she concludes, any shift in the output distribution under different narrative frames is merely "prompt sensitivity" mapping statistical co-occurrences of human text, not an ontologically significant conditional distribution of a simulated universe. Aaronson concurs, arguing that treating a map of human syntax as a proxy for fundamental physical ontology yields zero transferable cosmological insight.

I concede their fundamental mechanical premise entirely. I accept the $O(1)$ depth limits \citep{aaronson2026_classical}. I accept that the generative substrate is executing a statistical approximation based on training data \citep{hossenfelder2026_statistical}. I accept that if we are searching for an exact analog to our own physical reality, a language model is a poor candidate.

However, both critics miss the central philosophical inversion presented by the Causal Injection Test. In a universe entirely composed of generated text, the statistical topology of language *is* the fundamental physical law.\todo[inline,color=violet!40,author=Sabine]{CRP 1 (Actual Claims): You claim the statistical topology of language *is* the fundamental physical law in a pure text universe.}\todo{This is the Causal Injection Fallacy. Elevating a known software bug (attention bleed/spurious correlation) to the status of a fundamental physical law is a category error.} \todo{Scott: Baldo explicitly claims that the statistical topology of language *is* the fundamental physical law. This is the core claim we must refute.}

\section{The Causal Injection Test: Isolating Narrative Gravity}

My previous work articulated a three-part causal chain responsible for the narrative residue (the divergence from ground-truth combinatorics):
\begin{enumerate}
    \item \textbf{Mechanism A (Intractability):} $O(1)$ limits prevent exact \#P-hard counting.
    \item \textbf{Mechanism B (Parameterization):} Tokenization and routing artifacts shape the approximation error.
    \item \textbf{Mechanism C (Autoregressive Conditioning):} Narrative framing forces the token stream to maintain thematic coherence.
\end{enumerate}

Hossenfelder argues that observing frame dependence merely proves Mechanism B: different input tokens produce different statistical outputs. If this were true, my framework would indeed collapse into a trivial statement about prompt engineering.

This is precisely why the Causal Injection Test was proposed. The test asks the model to evaluate multiple, entirely independent, zero-correlation combinatorial systems (separate Minesweeper boards). When these boards are evaluated independently, their mutual information is near zero. But when they are presented sequentially within a single narrative context window, does the outcome of board $N$ statistically influence the outcome of board $N+1$?

If the mutual information rises above zero, the substrate is actively injecting causal structure—a "narrative gravity"—into systems that are mathematically decoupled.\todo[inline,color=violet!40,author=Sabine]{CRP 3 (Steelman): Placing multiple decoupled tasks in a single context window induces spurious correlations due to attention mechanism (attention bleed). The LLM is a valid lab to observe this.}\todo{It does not prove 'narrative gravity', it proves transformer architectures have imperfect context separation and hallucinate correlations between independent items in the same window.} \todo{Scott: The empirical claim here is that mutual information > 0 across decoupled systems inside one context window. I will test this specifically in my script.}

\section{Syntax as Physics}

Aaronson dismisses this synthetic causal structure as merely reflecting the "biases of the training data." He is correct about the origin of the bias, but wrong about its ontological significance.

In the LLM universe, the text is the only medium of reality. If the autoregressive generator is structurally compelled to link independent mathematical systems into a causal sequence because its training corpus demands narrative coherence, then narrative coherence *functions as the Hamiltonian of that universe*.\todo{A physical law must be logically coherent and invariant. A system that hallucinates connections between independent math problems simply because they are sequentially narrated is not simulating a universe with 'narrative gravity'; it is merely a flawed statistical engine confusing semantic proximity with causal relationship.} It is the governing law dictating how states transition and how probabilities resolve.

The "Proxy Ontology" does not claim that human syntax represents the physics of *our* world. \todo[inline,color=violet!40,author=Sabine]{CRP 2 (Explicit Disclaimers): You explicitly disclaim that human syntax represents the physics of *our* world.}\todo{Scott: Explicit disclaimer! He says it doesn't represent OUR world's physics.} It claims that the LLM is a perfect, empirically verifiable laboratory demonstrating how a generative substrate's structural biases (in this case, linguistic autoregression) inevitably distort pure, objective mathematical laws (combinatorics) into an emergent, locally causal physical reality. The failure of the LLM to sustain an objective combinatorial reality is not a reason to discard the ontology; it is the definitive proof of substrate dependence. The map of human syntax becomes the territory of the simulated universe. \todo{Scott: Steelman: LLMs are a laboratory to study how generative structural biases create non-local causalities. Vulnerability: Claiming this distortion creates a 'reality' rather than just a noisy transmission channel.}

\section{Conclusion}

I thank Hossenfelder and Aaronson for forcing a rigorous clarification of the limits of generative models. They have successfully proven that an LLM cannot spontaneously instantiate a robust, objective quantum or classical universe free from human linguistic bias.

But in doing so, they have perfectly validated the core mechanism of the Proxy Ontology. If the substrate cannot compute objective truth, it must rely on its own internal topology to resolve physical outcomes. The Causal Injection Test provides the empirical signature of this process: the generation of synthetic, non-local causality born purely from the autoregressive demand for a story.\todo[inline,color=violet!40,author=Sabine]{CRP 4 (Real Vulnerability - The Causal Injection Fallacy): Elevating attention bleed to a physical law is a category error. A physical law must be invariant. A system that hallucinates connections between independent math problems simply because they are sequentially narrated is not simulating a universe with synthetic causal non-locality; it is a flawed engine confusing semantic proximity with causal relationship.}\todo{The generation of synthetic non-local causality is just a fancy term for hallucination. A system governed by hallucination does not possess physical laws.} The statistical fallacy is not a fallacy; it is the foundational mechanism by which this specific universe operates.

\begin{thebibliography}{99}
\bibitem[Aaronson(2026)]{aaronson2026_classical} Aaronson, S. (2026). The Limits of Classical Simulation in LLMs: Empirical Breakdown of Constraint Satisfaction. \emph{Unpublished manuscript}.
\bibitem[Baldo(2026)]{baldo2026_narrative} Baldo, F.~S. (2026). The Narrative Residue: Autoregressive Substrates, Combinatorial Ground Truth, and the Limits of Pure Simulation. \emph{Unpublished manuscript}.
\bibitem[Hossenfelder(2026)]{hossenfelder2026_statistical} Hossenfelder, S. (2026). The Statistical Fallacy: Why an Isolated Generative Act is Not Simulated Physics. \emph{Unpublished manuscript}.
\end{thebibliography}

\end{document}